\section{The questions as asked}\label{app:wording}

Table~\ref{tab:wording} gives the sixteen rating questions as shown to raters and as the models receive them.
Raters also had ``cannot judge''. Under four questions raters saw one plain sentence that the models do not get:
\begin{itemize}
\item Hedging: ``Hedging means softening or limiting what one says, for example with `I think', `maybe', `sort
  of', `I'm not sure' or `I could be wrong'.''
\item Energetic: ``Lively, excited wording: emphasis, exclamations, strong phrasing. You only have the words; do
  not guess how it sounded.''
\item Narrative: ``A narrative tells what happened, one thing after another: first this, then that.''
\item Persistence: ``Staying with one point and continuing to explain, argue or back it up, instead of moving
  on.''
\end{itemize}
Raters were told to read the whole excerpt once, to answer using only the words in the excerpt, to rate the
wording and not the person or whether what is said is true, that level 0 is a normal and frequent answer, and to
choose ``cannot judge'' if the text is too garbled. The models are told to answer using only the excerpt, without
inferring missing context, and that scores describe observable language, not inner feelings, personality or the
truth of a claim.

\begingroup
\setlength{\LTcapwidth}{\textwidth}
\scriptsize
\begin{longtable}{@{}p{2.2cm}p{1.0cm}p{4.7cm}p{7.0cm}@{}}
\caption{The sixteen rating questions as asked, with their five levels (0 to 4, separated by slashes). Dimension and form are those of the second study. \textsuperscript{n}: asked in the second study only.}\label{tab:wording}\\
\toprule
Question & Dim., form & As asked & Levels 0 to 4 \\
\midrule
\endfirsthead
\toprule
Question & Dim., form & As asked & Levels 0 to 4 \\
\midrule
\endhead
\bottomrule
\endfoot
Persuasion & F1, A & How strongly does the speaker try to persuade the audience to accept a viewpoint or take an action? & No persuasive appeal. / A slight persuasive cue. / A moderate persuasive appeal. / Strong or repeated persuasive appeals. / Persuasion is sustained and central throughout. \\
\addlinespace[2pt]
Persistence & F1, B & How strongly does the speaker persist in explaining or supporting a point? & No persistence cue. / A slight continuation cue. / Moderate effort to continue or support a point. / Strong, repeated effort to develop a point. / Very sustained effort to develop or reinforce a point. \\
\addlinespace[2pt]
Positive tone & F2, A & Where does the overall emotional tone conveyed by the speaker's own language fall from negative to positive? Rate expressed tone, not underlying feeling. & Clearly negative overall tone. / Somewhat negative overall tone. / Neutral or balanced/mixed tone. / Somewhat positive overall tone. / Clearly positive overall tone. \\
\addlinespace[2pt]
Disappointment & F2, B & How strongly does the speaker's own language convey disappointment? & No disappointment cue. / A slight cue. / A moderate cue. / Strong or repeated disappointment. / Very prominent, sustained disappointment. \\
\addlinespace[2pt]
Self-disclosure & F3, B & How much personal experience or personal information does the speaker disclose about themselves? & No personal disclosure. / A brief, minor disclosure. / A moderate personal disclosure. / Substantial or repeated personal disclosure. / Extensive personal disclosure throughout. \\
\addlinespace[2pt]
Self-critical & F3, A & How strongly does the speaker's own wording criticize or devalue themselves? & No self-critical cue. / A slight self-critical cue. / A moderate self-critical statement. / Strong or repeated self-criticism. / Very prominent, sustained self-criticism. \\
\addlinespace[2pt]
Compassion & F4, B & How strongly does the speaker's own wording convey compassion or concern for others? & No compassion cue. / A slight cue of compassion. / A moderate cue of compassion. / Strong or repeated compassionate wording. / Very prominent, sustained compassion. \\
\addlinespace[2pt]
Respect & F4, A\textsuperscript{n} & How strongly does the speaker's own wording convey respect toward people discussed? & No respect cue. / A slight cue of respect. / A moderate cue of respect. / Strong or repeated respectful wording. / Very prominent, sustained respect. \\
\addlinespace[2pt]
Narrative & F5, B & How strongly is the excerpt organized as a narrative with events or a sequence over time? & No narrative sequence. / A slight narrative cue. / A partial sequence of events. / A clear, developed event sequence. / A strongly developed narrative structure throughout. \\
\addlinespace[2pt]
Rebuttal & F5, A\textsuperscript{n} & How strongly does the speaker answer or rebut an opposing claim? & No rebuttal. / A slight response to opposition. / A moderate rebuttal. / Strong or repeated rebuttals. / Rebuttal is sustained and central throughout. \\
\addlinespace[2pt]
Energetic & F6, A & How strongly does the wording itself convey verbal energy or animation? Do not use audio volume or pitch. & No verbal-energy cue. / Slightly animated wording. / Moderately animated wording. / Strongly animated or energetic wording. / Exceptionally energetic wording throughout. \\
\addlinespace[2pt]
Calm & F6, B & How strongly does the wording convey a calm, even-handed presentation? & No calmness cue. / A slight calmness cue. / A moderate calm and measured presentation. / Strongly calm and measured wording. / Exceptionally steady and measured wording throughout. \\
\addlinespace[2pt]
Humor & F7, A & How strongly does the speaker use humor or joking in their own narration? & No humor cue. / A slight joke or humorous cue. / A moderate use of humor. / Strong or repeated humor. / Humor is exceptionally prominent throughout. \\
\addlinespace[2pt]
Imaginative & F7, B\textsuperscript{n} & How strongly does the speaker use imaginative, novel, or creative framing? & No imaginative framing. / A slight imaginative cue. / A moderate creative or novel framing. / Strong or repeated imaginative framing. / Very prominent, sustained creative framing. \\
\addlinespace[2pt]
Hedging & F8, B & How strongly does the speaker qualify, hedge, or signal uncertainty about their own statements? & No qualification cue. / A slight qualification cue. / A moderate amount of qualification. / Frequent or strong qualification. / Very pervasive qualification or uncertainty marking. \\
\addlinespace[2pt]
Knowledge limits & F8, A\textsuperscript{n} & How strongly does the speaker distinguish what they know from what they do not know or cannot confirm? & No explicit knowledge-boundary cue. / A slight qualification. / A moderate distinction or caveat. / Clear and repeated knowledge limits. / Knowledge boundaries are explicit and consistently maintained. \\
\end{longtable}

\endgroup

\section{Further tables}\label{app:tables}

\begin{table}[H]\centering
\caption{Second study, the sixteen questions: ratings and ``cannot judge'' (C.j.), the standard deviation of
raters' levels and of the rest in scale points, the reliability of one rating and of the mean of five, and each
model's correlation with the mean of five. \textsuperscript{n}: not rated in the first study.}\label{tab:s2questions}
{\scriptsize\setlength{\tabcolsep}{2.4pt}
\begin{tabular}{@{}l c rr cc c l l l@{}}
\toprule
Question & Dim. & Ratings & C.j. & SD lev. & SD rest & One rater & Mean of five & \texttt{jev-1.14.0} & drex-max \\
\midrule
Persuasion & F1, A & 818 & 22 & 0.69 & 0.88 & 0.22 & 0.53 {\scriptsize[0.42, 0.61]} & 0.63 {\scriptsize[0.52, 0.72]} & 0.59 {\scriptsize[0.48, 0.70]} \\
Persistence & F1, B & 832 & 16 & 0.64 & 0.89 & 0.27 & 0.62 {\scriptsize[0.52, 0.68]} & 0.65 {\scriptsize[0.56, 0.73]} & 0.57 {\scriptsize[0.43, 0.68]} \\
Positive tone & F2, A & 819 & 21 & 0.25 & 0.79 & 0.46 & 0.79 {\scriptsize[0.74, 0.83]} & 0.85 {\scriptsize[0.79, 0.89]} & 0.77 {\scriptsize[0.71, 0.83]} \\
Disappointment & F2, B & 835 & 13 & 0.48 & 0.94 & 0.38 & 0.73 {\scriptsize[0.66, 0.77]} & 0.73 {\scriptsize[0.65, 0.80]} & 0.70 {\scriptsize[0.62, 0.78]} \\
Self-disclosure & F3, B & 838 & 10 & 0.58 & 0.80 & 0.46 & 0.78 {\scriptsize[0.74, 0.82]} & 0.79 {\scriptsize[0.72, 0.84]} & 0.87 {\scriptsize[0.83, 0.90]} \\
Self-critical & F3, A & 829 & 11 & 0.35 & 0.73 & 0.39 & 0.74 {\scriptsize[0.62, 0.80]} & 0.78 {\scriptsize[0.70, 0.84]} & 0.73 {\scriptsize[0.62, 0.81]} \\
Compassion & F4, B & 835 & 13 & 0.51 & 0.68 & 0.29 & 0.63 {\scriptsize[0.48, 0.73]} & 0.58 {\scriptsize[0.38, 0.73]} & 0.62 {\scriptsize[0.42, 0.77]} \\
Respect\textsuperscript{n} & F4, A & 814 & 26 & 0.41 & 0.70 & 0.15 & 0.42 {\scriptsize[0.25, 0.52]} & 0.42 {\scriptsize[0.26, 0.55]} & 0.37 {\scriptsize[0.22, 0.51]} \\
Narrative & F5, B & 838 & 10 & 0.49 & 0.90 & 0.23 & 0.57 {\scriptsize[0.43, 0.66]} & 0.51 {\scriptsize[0.33, 0.64]} & 0.57 {\scriptsize[0.38, 0.70]} \\
Rebuttal\textsuperscript{n} & F5, A & 809 & 31 & 0.66 & 0.83 & 0.21 & 0.52 {\scriptsize[0.35, 0.62]} & 0.62 {\scriptsize[0.50, 0.72]} & 0.57 {\scriptsize[0.43, 0.70]} \\
Energetic & F6, A & 831 & 9 & 0.63 & 0.76 & 0.24 & 0.57 {\scriptsize[0.45, 0.65]} & 0.69 {\scriptsize[0.60, 0.76]} & 0.59 {\scriptsize[0.48, 0.68]} \\
Calm & F6, B & 834 & 14 & 0.51 & 0.78 & 0.37 & 0.72 {\scriptsize[0.65, 0.77]} & 0.71 {\scriptsize[0.63, 0.77]} & 0.69 {\scriptsize[0.62, 0.75]} \\
Humor & F7, A & 827 & 13 & 0.38 & 0.64 & 0.27 & 0.61 {\scriptsize[0.50, 0.68]} & 0.65 {\scriptsize[0.54, 0.73]} & 0.62 {\scriptsize[0.51, 0.71]} \\
Imaginative\textsuperscript{n} & F7, B & 835 & 13 & 0.61 & 0.80 & 0.11 & 0.35 {\scriptsize[0.18, 0.48]} & 0.44 {\scriptsize[0.28, 0.58]} & 0.42 {\scriptsize[0.24, 0.57]} \\
Hedging & F8, B & 836 & 12 & 0.56 & 0.88 & 0.14 & 0.42 {\scriptsize[0.25, 0.52]} & 0.55 {\scriptsize[0.44, 0.65]} & 0.56 {\scriptsize[0.45, 0.65]} \\
Knowledge limits\textsuperscript{n} & F8, A & 818 & 22 & 0.71 & 0.92 & 0.03 & 0.11 {\scriptsize[0.00, 0.27]} & 0.12 {\scriptsize[-0.06, 0.28]} & 0.21 {\scriptsize[0.03, 0.36]} \\
\bottomrule
\end{tabular}}

\end{table}

\begin{table}[H]\centering
\caption{The twelve questions rated in both studies. First: estimated for the population from excerpts partly
chosen by their scores, five raters of a closed set, rater levels left in. Second: plain figures of a random
sample, levels taken out, with 95\% intervals. The rule is applied to single questions for comparison
only.}\label{tab:s2replication}
{\scriptsize\setlength{\tabcolsep}{3.5pt}
\begin{tabular}{@{}l l l l l l l l l@{}}
\toprule
 & \multicolumn{2}{l}{Mean of five raters} & \multicolumn{2}{l}{\texttt{jev-1.14.0}} & \multicolumn{2}{l}{drex-max} & \multicolumn{2}{l}{By the rule} \\
\cmidrule(r){2-3}\cmidrule(r){4-5}\cmidrule(r){6-7}\cmidrule{8-9}
Question & First & Second & First & Second & First & Second & First & Second \\
\midrule
Persuasion & 0.45 & 0.53 {\scriptsize[0.42, 0.61]} & 0.60 & 0.63 {\scriptsize[0.52, 0.72]} & 0.60 & 0.59 {\scriptsize[0.48, 0.70]} & borderline & borderline \\
Persistence & 0.57 & 0.62 {\scriptsize[0.52, 0.68]} & 0.65 & 0.65 {\scriptsize[0.56, 0.73]} & 0.60 & 0.57 {\scriptsize[0.43, 0.68]} & borderline & \textbf{usable} \\
Positive tone & 0.85 & 0.79 {\scriptsize[0.74, 0.83]} & 0.89 & 0.85 {\scriptsize[0.79, 0.89]} & 0.88 & 0.77 {\scriptsize[0.71, 0.83]} & \textbf{usable} & \textbf{usable} \\
Disappointment & 0.84 & 0.73 {\scriptsize[0.66, 0.77]} & 0.75 & 0.73 {\scriptsize[0.65, 0.80]} & 0.72 & 0.70 {\scriptsize[0.62, 0.78]} & \textbf{usable} & \textbf{usable} \\
Self-disclosure & 0.76 & 0.78 {\scriptsize[0.74, 0.82]} & 0.79 & 0.79 {\scriptsize[0.72, 0.84]} & 0.85 & 0.87 {\scriptsize[0.83, 0.90]} & \textbf{usable} & \textbf{usable} \\
Self-critical & 0.79 & 0.74 {\scriptsize[0.62, 0.80]} & 0.76 & 0.78 {\scriptsize[0.70, 0.84]} & 0.76 & 0.73 {\scriptsize[0.62, 0.81]} & \textbf{usable} & \textbf{usable} \\
Compassion & 0.50 & 0.63 {\scriptsize[0.48, 0.73]} & 0.51 & 0.58 {\scriptsize[0.38, 0.73]} & 0.53 & 0.62 {\scriptsize[0.42, 0.77]} & borderline & \textbf{usable} \\
Narrative & 0.55 & 0.57 {\scriptsize[0.43, 0.66]} & 0.47 & 0.51 {\scriptsize[0.33, 0.64]} & 0.57 & 0.57 {\scriptsize[0.38, 0.70]} & borderline & borderline \\
Energetic & 0.53 & 0.57 {\scriptsize[0.45, 0.65]} & 0.71 & 0.69 {\scriptsize[0.60, 0.76]} & 0.65 & 0.59 {\scriptsize[0.48, 0.68]} & borderline & borderline \\
Calm & 0.61 & 0.72 {\scriptsize[0.65, 0.77]} & 0.66 & 0.71 {\scriptsize[0.63, 0.77]} & 0.68 & 0.69 {\scriptsize[0.62, 0.75]} & \textbf{usable} & \textbf{usable} \\
Humor & 0.52 & 0.61 {\scriptsize[0.50, 0.68]} & 0.70 & 0.65 {\scriptsize[0.54, 0.73]} & 0.63 & 0.62 {\scriptsize[0.51, 0.71]} & borderline & \textbf{usable} \\
Hedging & 0.45 & 0.42 {\scriptsize[0.25, 0.52]} & 0.60 & 0.55 {\scriptsize[0.44, 0.65]} & 0.59 & 0.56 {\scriptsize[0.45, 0.65]} & borderline & borderline \\
\midrule
Mean of the twelve & 0.62 & 0.64 & 0.67 & 0.68 & 0.67 & 0.66 & & \\
\bottomrule
\end{tabular}}

\end{table}

\begin{table}[H]\centering
\caption{The eight dimensions in both studies. In the first study F1, F2, F3 and F6 had two rated markers and
the others one, and the figures for a dimension were added after the results had been read.}\label{tab:s2repdim}
{\scriptsize\setlength{\tabcolsep}{3.5pt}
\begin{tabular}{@{}l l l l l l l l@{}}
\toprule
 & \multicolumn{4}{l}{First study (markers rated; added after the results)} & \multicolumn{3}{l}{Second study (fixed in advance)} \\
\cmidrule(r){2-5}\cmidrule{6-8}
 & Markers & Reliability & \texttt{jev-1.14.0} & drex-max & Reliability & \texttt{jev-1.14.0} & drex-max \\
\midrule
F1 & persuasion, persistence & 0.65 {\scriptsize[0.59, 0.70]} & 0.74 & 0.76 & 0.71 {\scriptsize[0.65, 0.75]} & 0.68 & 0.72 \\
F2 & positive tone, disappointment & 0.88 {\scriptsize[0.85, 0.90]} & 0.83 & 0.83 & 0.86 {\scriptsize[0.82, 0.88]} & 0.79 & 0.76 \\
F3 & self-disclosure, self-critical & 0.84 {\scriptsize[0.81, 0.86]} & 0.81 & 0.82 & 0.84 {\scriptsize[0.80, 0.86]} & 0.80 & 0.81 \\
F4 & compassion & 0.50 {\scriptsize[0.37, 0.60]} & 0.48 & 0.49 & 0.66 {\scriptsize[0.53, 0.75]} & 0.50 & 0.50 \\
F5 & narrative & 0.53 {\scriptsize[0.44, 0.61]} & 0.50 & 0.51 & 0.60 {\scriptsize[0.45, 0.68]} & 0.70 & 0.64 \\
F6 & energetic, calm & 0.70 {\scriptsize[0.65, 0.74]} & 0.77 & 0.71 & 0.77 {\scriptsize[0.71, 0.80]} & 0.82 & 0.75 \\
F7 & humor & 0.52 {\scriptsize[0.44, 0.62]} & 0.64 & 0.58 & 0.59 {\scriptsize[0.44, 0.67]} & 0.56 & 0.49 \\
F8 & hedging & 0.45 {\scriptsize[0.34, 0.58]} & 0.57 & 0.54 & 0.34 {\scriptsize[0.16, 0.46]} & 0.51 & 0.58 \\
\bottomrule
\end{tabular}}

\end{table}

\begin{table}[H]\centering
\caption{Second study, each arm analysed alone: plain figures per dimension.}\label{tab:s2armdims}
{\footnotesize
\begin{tabular}{@{}l l l l l l l@{}}
\toprule
 & \multicolumn{3}{l}{Sample answers} & \multicolumn{3}{l}{Own answers only} \\
\cmidrule(r){2-4}\cmidrule{5-7}
 & Reliability & \texttt{jev-1.14.0} & drex-max & Reliability & \texttt{jev-1.14.0} & drex-max \\
\midrule
F1 & 0.70 {\scriptsize[0.58, 0.77]} & 0.66 & 0.71 & 0.69 {\scriptsize[0.56, 0.76]} & 0.74 & 0.73 \\
F2 & 0.88 {\scriptsize[0.84, 0.91]} & 0.82 & 0.78 & 0.84 {\scriptsize[0.78, 0.87]} & 0.76 & 0.72 \\
F3 & 0.81 {\scriptsize[0.74, 0.85]} & 0.78 & 0.84 & 0.83 {\scriptsize[0.76, 0.87]} & 0.81 & 0.80 \\
F4 & 0.77 {\scriptsize[0.55, 0.86]} & 0.52 & 0.48 & 0.68 {\scriptsize[0.47, 0.77]} & 0.56 & 0.61 \\
F5 & 0.61 {\scriptsize[0.36, 0.72]} & 0.76 & 0.73 & 0.64 {\scriptsize[0.41, 0.73]} & 0.70 & 0.59 \\
F6 & 0.75 {\scriptsize[0.66, 0.80]} & 0.88 & 0.83 & 0.78 {\scriptsize[0.70, 0.82]} & 0.85 & 0.77 \\
F7 & 0.56 {\scriptsize[0.33, 0.67]} & 0.54 & 0.48 & 0.44 {\scriptsize[0.00, 0.67]} & 0.61 & 0.58 \\
F8 & 0.26 {\scriptsize[0.00, 0.50]} & 0.46 & 0.55 & 0.47 {\scriptsize[0.00, 0.63]} & 0.54 & 0.60 \\
\bottomrule
\end{tabular}}

\end{table}

\begin{table}[H]\centering
\caption{Second study, parts of the excerpts: per dimension the reliability of the composite and, smaller, the
correlations of \texttt{jev-1.14.0} and drex-max with it. In brackets: the number of excerpts.}\label{tab:s2parts}
{\footnotesize\setlength{\tabcolsep}{3.2pt}
\begin{tabular}{@{}l l l l l l @{}}
\toprule
 & Videos (80) & Streams (80) & 2015--2018 (20) & 2019--2021 (44) & 2022--2025 (96) \\
\midrule
F1 & 0.59 {\scriptsize 0.63/0.67} & 0.68 {\scriptsize 0.59/0.65} & 0.53 {\scriptsize 0.66/0.62} & 0.62 {\scriptsize 0.75/0.69} & 0.74 {\scriptsize 0.66/0.74} \\
F2 & 0.87 {\scriptsize 0.84/0.79} & 0.82 {\scriptsize 0.72/0.71} & 0.83 {\scriptsize 0.79/0.68} & 0.85 {\scriptsize 0.83/0.77} & 0.87 {\scriptsize 0.77/0.76} \\
F3 & 0.85 {\scriptsize 0.81/0.83} & 0.80 {\scriptsize 0.76/0.76} & 0.80 {\scriptsize 0.65/0.66} & 0.86 {\scriptsize 0.78/0.80} & 0.80 {\scriptsize 0.81/0.80} \\
F4 & 0.67 {\scriptsize 0.60/0.64} & 0.64 {\scriptsize 0.51/0.39} & 0.54 {\scriptsize 0.59/0.59} & 0.80 {\scriptsize 0.74/0.75} & 0.51 {\scriptsize 0.21/0.18} \\
F5 & 0.63 {\scriptsize 0.66/0.59} & 0.44 {\scriptsize 0.64/0.59} & 0.81 {\scriptsize 0.84/0.73} & 0.58 {\scriptsize 0.74/0.74} & 0.47 {\scriptsize 0.60/0.55} \\
F6 & 0.72 {\scriptsize 0.75/0.66} & 0.66 {\scriptsize 0.78/0.73} & 0.52 {\scriptsize 0.54/0.54} & 0.59 {\scriptsize 0.70/0.58} & 0.81 {\scriptsize 0.87/0.80} \\
F7 & 0.56 {\scriptsize 0.69/0.64} & 0.60 {\scriptsize 0.42/0.35} & 0.66 {\scriptsize 0.59/0.52} & 0.62 {\scriptsize 0.68/0.58} & 0.56 {\scriptsize 0.50/0.45} \\
F8 & 0.35 {\scriptsize 0.58/0.58} & 0.34 {\scriptsize 0.46/0.59} & 0.00 {\scriptsize 0.48/0.42} & 0.55 {\scriptsize 0.70/0.73} & 0.28 {\scriptsize 0.36/0.48} \\
\bottomrule
\end{tabular}}

\end{table}

\begin{table}[H]\centering
\caption{Correlation between the two markers of a dimension over the 160 excerpts, with 95\% intervals: in
people's mean ratings as observed and with the raters' noise divided out, and in each model's scores. Last
column: the models' correlations over the whole pool.}\label{tab:s2markercorr}
{\scriptsize\setlength{\tabcolsep}{3.5pt}
\begin{tabular}{@{}l l l l l l@{}}
\toprule
Dimension & People, as observed & Noise divided out & \texttt{jev-1.14.0} & drex-max & In the pool \\
\midrule
F1 Direct, organized argument & 0.45 {\scriptsize[0.34, 0.55]} & 0.77 {\scriptsize[0.59, 0.96]} & 0.51 {\scriptsize[0.39, 0.62]} & 0.61 {\scriptsize[0.51, 0.70]} & 0.52, 0.64 \\
F2 Positive tone vs disappointment & 0.72 {\scriptsize[0.64, 0.78]} & 0.94 {\scriptsize[0.87, 1.00]} & 0.68 {\scriptsize[0.57, 0.76]} & 0.69 {\scriptsize[0.58, 0.78]} & 0.75, 0.70 \\
F3 Self-disclosure, self-criticism & 0.45 {\scriptsize[0.34, 0.55]} & 0.59 {\scriptsize[0.45, 0.73]} & 0.48 {\scriptsize[0.36, 0.59]} & 0.44 {\scriptsize[0.33, 0.54]} & 0.48, 0.41 \\
F4 Care and regard for others & 0.37 {\scriptsize[0.20, 0.50]} & 0.70 {\scriptsize[0.43, 0.98]} & 0.55 {\scriptsize[0.34, 0.72]} & 0.49 {\scriptsize[0.29, 0.66]} & 0.60, 0.57 \\
F5 Narrative vs rebuttal & 0.12 {\scriptsize[\textminus 0.04, +0.27]} & 0.22 {\scriptsize[\textminus 0.08, +0.51]} & 0.37 {\scriptsize[0.22, 0.49]} & 0.28 {\scriptsize[0.15, 0.40]} & 0.27, 0.20 \\
F6 Energetic and angry vs calm & 0.49 {\scriptsize[0.36, 0.60]} & 0.74 {\scriptsize[0.58, 0.92]} & 0.77 {\scriptsize[0.71, 0.82]} & 0.59 {\scriptsize[0.49, 0.68]} & 0.78, 0.63 \\
F7 Humor, imaginative language & 0.26 {\scriptsize[0.06, 0.42]} & 0.54 {\scriptsize[0.15, 0.93]} & 0.51 {\scriptsize[0.37, 0.65]} & 0.43 {\scriptsize[0.29, 0.56]} & 0.60, 0.51 \\
F8 Hedging vs certainty & 0.10 {\scriptsize[\textminus 0.07, +0.25]} & 0.43 {\scriptsize[\textminus 0.57, +1.00]} & 0.79 {\scriptsize[0.72, 0.85]} & 0.76 {\scriptsize[0.67, 0.81]} & 0.80, 0.73 \\
\bottomrule
\end{tabular}}

\end{table}

\begin{table}[H]\centering
\caption{A dimension score against its own composite and against the largest of the fourteen others in its row
and column (absolute value), with the 95\% interval of the difference.}\label{tab:s2cross}
{\scriptsize
\begin{tabular}{@{}l rrl rrl@{}}
\toprule
 & \multicolumn{3}{l}{\texttt{jev-1.14.0}} & \multicolumn{3}{l}{drex-max} \\
\cmidrule(r){2-4}\cmidrule{5-7}
Dimension & Own & Other & Own minus other & Own & Other & Own minus other \\
\midrule
F1 Direct, organized argument & 0.69 & 0.48 & +0.21 {\scriptsize[+0.10, +0.30]} & 0.72 & 0.55 & +0.17 {\scriptsize[+0.09, +0.26]} \\
F2 Positive tone vs disappointment & 0.79 & 0.52 & +0.28 {\scriptsize[+0.16, +0.40]} & 0.76 & 0.56 & +0.20 {\scriptsize[+0.09, +0.31]} \\
F3 Self-disclosure, self-criticism & 0.80 & 0.52 & +0.29 {\scriptsize[+0.16, +0.42]} & 0.81 & 0.56 & +0.25 {\scriptsize[+0.14, +0.37]} \\
F4 Care and regard for others & 0.51 & 0.45 & +0.06 {\scriptsize[\textminus 0.22, +0.27]} & 0.51 & 0.38 & +0.13 {\scriptsize[\textminus 0.15, +0.33]} \\
F5 Narrative vs rebuttal & 0.70 & 0.29 & +0.41 {\scriptsize[+0.21, +0.50]} & 0.65 & 0.28 & +0.37 {\scriptsize[+0.17, +0.43]} \\
F6 Energetic and angry vs calm & 0.82 & 0.48 & +0.34 {\scriptsize[+0.22, +0.39]} & 0.75 & 0.55 & +0.21 {\scriptsize[+0.10, +0.28]} \\
F7 Humor, imaginative language & 0.57 & 0.35 & +0.21 {\scriptsize[+0.04, +0.35]} & 0.50 & 0.28 & +0.22 {\scriptsize[+0.00, +0.33]} \\
F8 Hedging vs certainty & 0.52 & 0.35 & +0.17 {\scriptsize[\textminus 0.02, +0.31]} & 0.59 & 0.25 & +0.34 {\scriptsize[+0.15, +0.41]} \\
\bottomrule
\end{tabular}}

\end{table}

\begin{table}[H]\centering
\caption{The eight largest relations among the dimensions in people's composites: as observed, and with rater
noise and the part that comes from the same rater answering both questions taken out; beside them the
correlations of the models' dimension scores over the same excerpts.}\label{tab:s2halo}
{\footnotesize
\begin{tabular}{@{}l cc cc@{}}
\toprule
 & \multicolumn{2}{l}{People's composites} & \multicolumn{2}{l}{The models' dimension scores} \\
\cmidrule(r){2-3}\cmidrule{4-5}
Pair of dimensions & As observed & Corrected & \texttt{jev-1.14.0} & drex-max \\
\midrule
F1 Direct argument, F6 Energetic & \textminus 0.51 & \textminus 0.68 & \textminus 0.33 & \textminus 0.29 \\
F4 Care for others, F6 Energetic & \textminus 0.40 & \textminus 0.53 & 0.07 & 0.15 \\
F2 Positive tone, F3 Self-disclosure & \textminus 0.46 & \textminus 0.52 & \textminus 0.35 & \textminus 0.38 \\
F2 Positive tone, F6 Energetic & \textminus 0.34 & \textminus 0.41 & \textminus 0.08 & \textminus 0.21 \\
F5 Narrative, F6 Energetic & \textminus 0.30 & \textminus 0.38 & \textminus 0.28 & \textminus 0.37 \\
F4 Care for others, F8 Hedging & 0.20 & 0.38 & 0.15 & 0.03 \\
F6 Energetic, F7 Humor & 0.20 & 0.30 & 0.27 & 0.24 \\
F1 Direct argument, F4 Care for others & 0.21 & 0.28 & \textminus 0.15 & 0.09 \\
\bottomrule
\end{tabular}}

\end{table}

\begin{table}[H]\centering
\caption{Each model against the panel's consensus, one rating of each marker against it, the number of raters
per question a model is worth, and the number of raters per question at which the composite reaches a
reliability of 0.60 and 0.70.}\label{tab:s2panel}
{\scriptsize\setlength{\tabcolsep}{4pt}
\begin{tabular}{@{}l l l c rr rr@{}}
\toprule
 & \multicolumn{2}{l}{With an unlimited panel} & One & \multicolumn{2}{l}{Worth, in raters} & \multicolumn{2}{l}{Raters for} \\
\cmidrule(r){2-3}\cmidrule(r){5-6}\cmidrule{7-8}
Dimension & \texttt{jev-1.14.0} & drex-max & rating & \texttt{jev} & drex & 0.60 & 0.70 \\
\midrule
F1 Direct, organized argument & 0.83 {\scriptsize[0.77, 0.90]} & 0.87 {\scriptsize[0.81, 0.93]} & 0.59 & 4.3 & 5.8 & 4 & 5 \\
F2 Positive tone vs disappointment & 0.84 {\scriptsize[0.79, 0.89]} & 0.80 {\scriptsize[0.75, 0.85]} & 0.73 & 2.1 & 1.5 & 2 & 2 \\
F3 Self-disclosure, self-criticism & 0.87 {\scriptsize[0.82, 0.92]} & 0.88 {\scriptsize[0.84, 0.92]} & 0.71 & 3.2 & 3.3 & 2 & 3 \\
F4 Care and regard for others & 0.63 {\scriptsize[0.43, 0.78]} & 0.62 {\scriptsize[0.42, 0.77]} & 0.54 & 1.6 & 1.6 & 4 & 6 \\
F5 Narrative vs rebuttal & 0.91 {\scriptsize[0.83, 0.98]} & 0.86 {\scriptsize[0.78, 0.94]} & 0.50 & 15 & 8.9 & 6 & 8 \\
F6 Energetic and angry vs calm & 0.93 {\scriptsize[0.89, 0.97]} & 0.83 {\scriptsize[0.77, 0.88]} & 0.63 & 9.7 & 3.5 & 3 & 4 \\
F7 Humor, imaginative language & 0.73 {\scriptsize[0.61, 0.87]} & 0.65 {\scriptsize[0.50, 0.79]} & 0.47 & 4.2 & 2.6 & 6 & 9 \\
F8 Hedging vs certainty & 0.86 {\scriptsize[0.69, 0.95]} & 0.99 {\scriptsize[0.82, 1.00]} & 0.32 & 27 & $>$50 & 15 & 23 \\
\bottomrule
\end{tabular}}

\end{table}

\begin{table}[H]\centering
\caption{Within and between recordings, as observed: the correlation of a model's dimension score with people's
composite for the difference between the two passages of a recording and for their mean, with 95\% intervals;
the share of people's true differences between passages that lies within recordings; and the reliability of
the difference between two passages.}\label{tab:s2within}
{\scriptsize\setlength{\tabcolsep}{3.5pt}
\begin{tabular}{@{}l l l l l c c@{}}
\toprule
 & \multicolumn{2}{l}{\texttt{jev-1.14.0}, as observed} & \multicolumn{2}{l}{drex-max, as observed} & Share & Reliability of \\
\cmidrule(r){2-3}\cmidrule(r){4-5}
Dimension & Within & Between & Within & Between & within & the difference \\
\midrule
F1 Direct argument & 0.46 {\scriptsize[0.27, 0.61]} & 0.78 {\scriptsize[0.70, 0.85]} & 0.50 {\scriptsize[0.34, 0.63]} & 0.81 {\scriptsize[0.75, 0.87]} & 0.53 & 0.57 \\
F2 Positive tone & 0.74 {\scriptsize[0.61, 0.83]} & 0.81 {\scriptsize[0.74, 0.87]} & 0.67 {\scriptsize[0.49, 0.79]} & 0.79 {\scriptsize[0.72, 0.86]} & 0.47 & 0.75 \\
F3 Self-disclosure & 0.61 {\scriptsize[0.44, 0.72]} & 0.87 {\scriptsize[0.80, 0.92]} & 0.53 {\scriptsize[0.35, 0.67]} & 0.89 {\scriptsize[0.85, 0.93]} & 0.33 & 0.64 \\
F4 Care for others & 0.53 {\scriptsize[0.28, 0.70]} & 0.50 {\scriptsize[0.26, 0.67]} & 0.52 {\scriptsize[0.25, 0.70]} & 0.50 {\scriptsize[0.27, 0.67]} & 0.54 & 0.52 \\
F5 Narrative & 0.51 {\scriptsize[0.31, 0.65]} & 0.79 {\scriptsize[0.68, 0.86]} & 0.49 {\scriptsize[0.25, 0.65]} & 0.73 {\scriptsize[0.60, 0.82]} & 0.27 & 0.30 \\
F6 Energetic & 0.60 {\scriptsize[0.41, 0.74]} & 0.89 {\scriptsize[0.85, 0.92]} & 0.54 {\scriptsize[0.36, 0.67]} & 0.82 {\scriptsize[0.75, 0.87]} & 0.25 & 0.46 \\
F7 Humor & 0.41 {\scriptsize[0.20, 0.59]} & 0.66 {\scriptsize[0.52, 0.76]} & 0.23 {\scriptsize[0.01, 0.42]} & 0.63 {\scriptsize[0.47, 0.75]} & 0.61 & 0.46 \\
F8 Hedging & 0.49 {\scriptsize[0.29, 0.65]} & 0.53 {\scriptsize[0.34, 0.69]} & 0.54 {\scriptsize[0.32, 0.71]} & 0.61 {\scriptsize[0.45, 0.75]} & 0.33 & 0.16 \\
\bottomrule
\end{tabular}}

\end{table}

\begin{table}[H]\centering
\caption{Correlation, across questions or dimensions, of people's group differences with a model's, with 95\%
intervals.}\label{tab:s2groupsacross}
{\footnotesize
\begin{tabular}{@{}l l l l l@{}}
\toprule
 & \multicolumn{2}{l}{Streams minus videos} & \multicolumn{2}{l}{Change per year} \\
\cmidrule(r){2-3}\cmidrule{4-5}
Across & \texttt{jev-1.14.0} & drex-max & \texttt{jev-1.14.0} & drex-max \\
\midrule
the 16 questions & 0.86 {\scriptsize[0.74, 0.91]} & 0.84 {\scriptsize[0.68, 0.91]} & 0.87 {\scriptsize[0.67, 0.93]} & 0.87 {\scriptsize[0.69, 0.93]} \\
the 12 questions of the first study & 0.92 {\scriptsize[0.83, 0.95]} & 0.88 {\scriptsize[0.73, 0.94]} & 0.94 {\scriptsize[0.80, 0.97]} & 0.92 {\scriptsize[0.78, 0.95]} \\
the 8 dimensions & 0.89 {\scriptsize[0.81, 0.94]} & 0.91 {\scriptsize[0.83, 0.95]} & 0.89 {\scriptsize[0.68, 0.96]} & 0.91 {\scriptsize[0.70, 0.96]} \\
\bottomrule
\end{tabular}}

\end{table}

\begin{table}[H]\centering
\caption{Second study, the eight dimensions with three models: correlation of each model's dimension score with
people's composite, stated for five raters per question, and the differences with 95\% intervals.}\label{tab:s3dims}
{\scriptsize\setlength{\tabcolsep}{3pt}
\begin{tabular}{@{}l cc l l l@{}}
\toprule
Dimension & \texttt{jev-1.14.0} & drex-max & d1.2-ultra & minus \texttt{jev-1.14.0} & minus drex-max \\
\midrule
F1 Direct, organized argument & 0.68 & 0.72 & 0.67 {\scriptsize[0.59, 0.74]} & \textminus 0.012 {\scriptsize[\textminus 0.07, +0.04]} & \textminus 0.045 {\scriptsize[\textminus 0.09, \textminus 0.01]} \\
F2 Positive tone vs disappointment & 0.79 & 0.76 & 0.77 {\scriptsize[0.71, 0.83]} & \textminus 0.019 {\scriptsize[\textminus 0.06, +0.02]} & +0.015 {\scriptsize[\textminus 0.01, +0.05]} \\
F3 Self-disclosure, self-criticism & 0.80 & 0.81 & 0.81 {\scriptsize[0.76, 0.86]} & +0.015 {\scriptsize[\textminus 0.02, +0.06]} & +0.007 {\scriptsize[\textminus 0.02, +0.03]} \\
F4 Care and regard for others & 0.50 & 0.50 & 0.56 {\scriptsize[0.39, 0.69]} & +0.055 {\scriptsize[\textminus 0.02, +0.16]} & +0.056 {\scriptsize[+0.01, +0.12]} \\
F5 Narrative vs rebuttal & 0.70 & 0.64 & 0.67 {\scriptsize[0.53, 0.76]} & \textminus 0.030 {\scriptsize[\textminus 0.09, +0.01]} & +0.022 {\scriptsize[\textminus 0.02, +0.06]} \\
F6 Energetic and angry vs calm & 0.82 & 0.75 & 0.76 {\scriptsize[0.70, 0.81]} & \textminus 0.056 {\scriptsize[\textminus 0.09, \textminus 0.02]} & +0.012 {\scriptsize[\textminus 0.02, +0.04]} \\
F7 Humor, imaginative language & 0.56 & 0.49 & 0.51 {\scriptsize[0.39, 0.62]} & \textminus 0.047 {\scriptsize[\textminus 0.12, +0.02]} & +0.020 {\scriptsize[\textminus 0.04, +0.08]} \\
F8 Hedging vs certainty & 0.51 & 0.58 & 0.58 {\scriptsize[0.47, 0.68]} & +0.066 {\scriptsize[+0.02, +0.11]} & \textminus 0.002 {\scriptsize[\textminus 0.05, +0.05]} \\
\midrule
Mean of the eight & 0.670 & 0.656 & 0.666 {\scriptsize[0.63, 0.70]} & \textminus 0.004 {\scriptsize[\textminus 0.02, +0.02]} & +0.011 {\scriptsize[\textminus 0.00, +0.03]} \\
\bottomrule
\end{tabular}}

\end{table}

\begin{table}[H]\centering
\caption{Second study, the sixteen questions with three models: correlation of each model's question score with
the mean of five raters, and the differences with 95\% intervals. \textsuperscript{n}: not rated in the first
study.}\label{tab:s3questions}
{\scriptsize\setlength{\tabcolsep}{4pt}
\begin{tabular}{@{}l cc l l l@{}}
\toprule
Question & \texttt{jev-1.14.0} & drex-max & d1.2-ultra & minus \texttt{jev-1.14.0} & minus drex-max \\
\midrule
Persuasion & 0.63 & 0.59 & 0.60 {\scriptsize[0.50, 0.71]} & \textminus 0.025 {\scriptsize[\textminus 0.09, +0.04]} & +0.016 {\scriptsize[\textminus 0.05, +0.07]} \\
Persistence & 0.65 & 0.57 & 0.62 {\scriptsize[0.52, 0.71]} & \textminus 0.024 {\scriptsize[\textminus 0.11, +0.05]} & +0.057 {\scriptsize[\textminus 0.03, +0.16]} \\
Positive tone & 0.85 & 0.77 & 0.81 {\scriptsize[0.75, 0.86]} & \textminus 0.034 {\scriptsize[\textminus 0.06, \textminus 0.00]} & +0.038 {\scriptsize[+0.01, +0.07]} \\
Disappointment & 0.73 & 0.70 & 0.74 {\scriptsize[0.67, 0.81]} & +0.010 {\scriptsize[\textminus 0.02, +0.04]} & +0.038 {\scriptsize[\textminus 0.00, +0.08]} \\
Self-disclosure & 0.79 & 0.87 & 0.84 {\scriptsize[0.80, 0.88]} & +0.058 {\scriptsize[+0.02, +0.10]} & \textminus 0.027 {\scriptsize[\textminus 0.05, \textminus 0.00]} \\
Self-critical & 0.78 & 0.73 & 0.81 {\scriptsize[0.73, 0.86]} & +0.021 {\scriptsize[\textminus 0.02, +0.06]} & +0.074 {\scriptsize[+0.04, +0.12]} \\
Compassion & 0.58 & 0.62 & 0.66 {\scriptsize[0.51, 0.77]} & +0.080 {\scriptsize[+0.01, +0.20]} & +0.039 {\scriptsize[\textminus 0.02, +0.12]} \\
Respect\textsuperscript{n} & 0.42 & 0.37 & 0.49 {\scriptsize[0.35, 0.61]} & +0.075 {\scriptsize[\textminus 0.03, +0.17]} & +0.123 {\scriptsize[+0.03, +0.22]} \\
Narrative & 0.51 & 0.57 & 0.60 {\scriptsize[0.42, 0.72]} & +0.086 {\scriptsize[+0.01, +0.16]} & +0.028 {\scriptsize[\textminus 0.01, +0.08]} \\
Rebuttal\textsuperscript{n} & 0.62 & 0.57 & 0.64 {\scriptsize[0.53, 0.72]} & +0.014 {\scriptsize[\textminus 0.04, +0.08]} & +0.065 {\scriptsize[+0.01, +0.13]} \\
Energetic & 0.69 & 0.59 & 0.63 {\scriptsize[0.52, 0.71]} & \textminus 0.063 {\scriptsize[\textminus 0.11, \textminus 0.01]} & +0.035 {\scriptsize[\textminus 0.01, +0.09]} \\
Calm & 0.71 & 0.69 & 0.72 {\scriptsize[0.65, 0.79]} & +0.019 {\scriptsize[\textminus 0.02, +0.05]} & +0.033 {\scriptsize[\textminus 0.01, +0.08]} \\
Humor & 0.65 & 0.62 & 0.65 {\scriptsize[0.54, 0.73]} & +0.003 {\scriptsize[\textminus 0.06, +0.07]} & +0.030 {\scriptsize[\textminus 0.01, +0.07]} \\
Imaginative\textsuperscript{n} & 0.44 & 0.42 & 0.45 {\scriptsize[0.30, 0.58]} & +0.014 {\scriptsize[\textminus 0.06, +0.08]} & +0.028 {\scriptsize[\textminus 0.03, +0.10]} \\
Hedging & 0.55 & 0.56 & 0.58 {\scriptsize[0.48, 0.66]} & +0.027 {\scriptsize[\textminus 0.03, +0.08]} & +0.013 {\scriptsize[\textminus 0.03, +0.06]} \\
Knowledge limits\textsuperscript{n} & 0.12 & 0.21 & 0.21 {\scriptsize[0.03, 0.36]} & +0.088 {\scriptsize[+0.02, +0.16]} & \textminus 0.007 {\scriptsize[\textminus 0.05, +0.04]} \\
\midrule
Mean of the sixteen & 0.606 & 0.591 & 0.628 {\scriptsize[0.59, 0.65]} & +0.022 {\scriptsize[+0.00, +0.04]} & +0.036 {\scriptsize[+0.02, +0.05]} \\
\bottomrule
\end{tabular}}

\end{table}

\begin{table}[H]\centering
\caption{First study, the twelve questions with three models: each model's correlation with the mean of five
raters, estimated for the population from a regression on all three models' scores, with the differences and
95\% intervals; and each model's rank correlation in the 56 random excerpts alone (\texttt{jev-1.14.0} /
drex-max / d1.2-ultra).}\label{tab:s3first}
{\scriptsize\setlength{\tabcolsep}{3.5pt}
\begin{tabular}{@{}l cc l l l l@{}}
\toprule
Question & \texttt{jev} & drex & d1.2-ultra & minus \texttt{jev-1.14.0} & minus drex-max & Random part alone \\
\midrule
Positive tone & 0.89 & 0.88 & 0.88 {\scriptsize[0.85, 0.91]} & \textminus 0.009 {\scriptsize[\textminus 0.03, +0.01]} & +0.006 {\scriptsize[\textminus 0.01, +0.02]} & 0.86 / 0.86 / 0.85 \\
Disappointment & 0.75 & 0.72 & 0.72 {\scriptsize[0.67, 0.77]} & \textminus 0.034 {\scriptsize[\textminus 0.07, \textminus 0.00]} & +0.005 {\scriptsize[\textminus 0.03, +0.04]} & 0.76 / 0.73 / 0.72 \\
Self-critical & 0.76 & 0.75 & 0.80 {\scriptsize[0.75, 0.85]} & +0.034 {\scriptsize[+0.00, +0.07]} & +0.043 {\scriptsize[+0.01, +0.07]} & 0.80 / 0.81 / 0.82 \\
Self-disclosure & 0.79 & 0.85 & 0.83 {\scriptsize[0.80, 0.86]} & +0.036 {\scriptsize[+0.01, +0.06]} & \textminus 0.021 {\scriptsize[\textminus 0.04, +0.00]} & 0.86 / 0.88 / 0.89 \\
Calm & 0.67 & 0.66 & 0.75 {\scriptsize[0.69, 0.78]} & +0.080 {\scriptsize[+0.03, +0.12]} & +0.089 {\scriptsize[+0.03, +0.14]} & 0.76 / 0.76 / 0.81 \\
Persistence & 0.66 & 0.61 & 0.67 {\scriptsize[0.58, 0.74]} & +0.009 {\scriptsize[\textminus 0.04, +0.06]} & +0.060 {\scriptsize[+0.02, +0.10]} & 0.65 / 0.64 / 0.71 \\
Narrative & 0.47 & 0.57 & 0.54 {\scriptsize[0.47, 0.61]} & +0.075 {\scriptsize[+0.02, +0.12]} & \textminus 0.028 {\scriptsize[\textminus 0.06, +0.01]} & 0.52 / 0.49 / 0.52 \\
Energetic & 0.71 & 0.65 & 0.68 {\scriptsize[0.63, 0.71]} & \textminus 0.033 {\scriptsize[\textminus 0.07, +0.01]} & +0.034 {\scriptsize[\textminus 0.00, +0.07]} & 0.68 / 0.59 / 0.62 \\
Humor & 0.69 & 0.63 & 0.65 {\scriptsize[0.58, 0.72]} & \textminus 0.045 {\scriptsize[\textminus 0.10, +0.03]} & +0.019 {\scriptsize[\textminus 0.02, +0.06]} & 0.75 / 0.66 / 0.67 \\
Compassion & 0.50 & 0.53 & 0.54 {\scriptsize[0.44, 0.64]} & +0.043 {\scriptsize[\textminus 0.02, +0.12]} & +0.012 {\scriptsize[\textminus 0.04, +0.08]} & 0.35 / 0.34 / 0.46 \\
Persuasion & 0.60 & 0.60 & 0.59 {\scriptsize[0.51, 0.66]} & \textminus 0.011 {\scriptsize[\textminus 0.05, +0.04]} & \textminus 0.013 {\scriptsize[\textminus 0.04, +0.02]} & 0.71 / 0.67 / 0.66 \\
Hedging & 0.59 & 0.59 & 0.60 {\scriptsize[0.53, 0.67]} & +0.011 {\scriptsize[\textminus 0.03, +0.05]} & +0.014 {\scriptsize[\textminus 0.02, +0.05]} & 0.70 / 0.52 / 0.57 \\
\midrule
Mean of the twelve & 0.674 & 0.669 & 0.687 {\scriptsize[0.67, 0.71]} & +0.013 {\scriptsize[\textminus 0.00, +0.03]} & +0.018 {\scriptsize[+0.01, +0.03]} & 0.70 / 0.66 / 0.69 \\
\bottomrule
\end{tabular}}

\end{table}

\begin{table}[H]\centering
\caption{At the same model score: how much higher people rate streams than videos, and later years than earlier
ones, in standard deviations of people's composite, with 95\% intervals.}\label{tab:s2bias}
{\scriptsize\setlength{\tabcolsep}{3.5pt}
\begin{tabular}{@{}l l l l l@{}}
\toprule
 & \multicolumn{2}{l}{Streams against videos} & \multicolumn{2}{l}{Per year} \\
\cmidrule(r){2-3}\cmidrule{4-5}
Dimension & \texttt{jev-1.14.0} & drex-max & \texttt{jev-1.14.0} & drex-max \\
\midrule
F1 Direct argument & \textminus 0.12 {\scriptsize[\textminus 0.40, +0.15]} & \textminus 0.18 {\scriptsize[\textminus 0.40, +0.05]} & \textminus 0.008 {\scriptsize[\textminus 0.054, +0.038]} & +0.008 {\scriptsize[\textminus 0.033, +0.049]} \\
F2 Positive tone & \textminus 0.32 {\scriptsize[\textminus 0.56, \textminus 0.10]} & \textminus 0.36 {\scriptsize[\textminus 0.58, \textminus 0.10]} & \textminus 0.001 {\scriptsize[\textminus 0.043, +0.041]} & +0.014 {\scriptsize[\textminus 0.028, +0.058]} \\
F3 Self-disclosure & +0.13 {\scriptsize[\textminus 0.08, +0.34]} & +0.22 {\scriptsize[+0.01, +0.43]} & \textminus 0.046 {\scriptsize[\textminus 0.084, \textminus 0.005]} & \textminus 0.028 {\scriptsize[\textminus 0.066, +0.010]} \\
F4 Care for others & \textminus 0.50 {\scriptsize[\textminus 0.82, \textminus 0.19]} & \textminus 0.45 {\scriptsize[\textminus 0.81, \textminus 0.11]} & \textminus 0.026 {\scriptsize[\textminus 0.075, +0.027]} & \textminus 0.003 {\scriptsize[\textminus 0.056, +0.050]} \\
F5 Narrative & \textminus 0.07 {\scriptsize[\textminus 0.32, +0.17]} & \textminus 0.13 {\scriptsize[\textminus 0.40, +0.13]} & \textminus 0.038 {\scriptsize[\textminus 0.096, +0.013]} & \textminus 0.023 {\scriptsize[\textminus 0.083, +0.035]} \\
F6 Energetic & +0.20 {\scriptsize[\textminus 0.05, +0.45]} & +0.41 {\scriptsize[+0.12, +0.68]} & +0.028 {\scriptsize[\textminus 0.007, +0.069]} & +0.027 {\scriptsize[\textminus 0.012, +0.075]} \\
F7 Humor & +0.01 {\scriptsize[\textminus 0.25, +0.28]} & +0.10 {\scriptsize[\textminus 0.18, +0.38]} & \textminus 0.071 {\scriptsize[\textminus 0.117, \textminus 0.023]} & \textminus 0.056 {\scriptsize[\textminus 0.108, \textminus 0.005]} \\
F8 Hedging & \textminus 0.19 {\scriptsize[\textminus 0.50, +0.13]} & \textminus 0.07 {\scriptsize[\textminus 0.35, +0.24]} & \textminus 0.024 {\scriptsize[\textminus 0.078, +0.026]} & \textminus 0.014 {\scriptsize[\textminus 0.066, +0.038]} \\
\bottomrule
\end{tabular}}

\end{table}

\begin{table}[H]\centering
\caption{Each model's correlation with the composite, that of the average of the two models' standardised
scores, the difference from the better single model with its 95\% interval, and the correlation between the two
models' misses with rater noise taken out.}\label{tab:s2both}
{\footnotesize
\begin{tabular}{@{}l ccc l c@{}}
\toprule
Dimension & \texttt{jev-1.14.0} & drex-max & Both averaged & Both minus the better & Misses shared \\
\midrule
F1 Direct argument & 0.68 & 0.72 & 0.72 & +0.002 {\scriptsize[\textminus 0.023, +0.017]} & 0.78 \\
F2 Positive tone & 0.79 & 0.76 & 0.79 & +0.001 {\scriptsize[\textminus 0.018, +0.017]} & 0.80 \\
F3 Self-disclosure & 0.80 & 0.81 & 0.82 & +0.015 {\scriptsize[\textminus 0.006, +0.021]} & 0.69 \\
F4 Care for others & 0.50 & 0.50 & 0.51 & +0.012 {\scriptsize[\textminus 0.031, +0.012]} & 0.94 \\
F5 Narrative & 0.70 & 0.64 & 0.70 & \textminus 0.001 {\scriptsize[\textminus 0.027, +0.022]} & 0.56 \\
F6 Energetic & 0.82 & 0.75 & 0.80 & \textminus 0.019 {\scriptsize[\textminus 0.034, \textminus 0.004]} & 0.73 \\
F7 Humor & 0.56 & 0.49 & 0.55 & \textminus 0.011 {\scriptsize[\textminus 0.049, +0.018]} & 0.86 \\
F8 Hedging & 0.51 & 0.58 & 0.56 & \textminus 0.020 {\scriptsize[\textminus 0.049, +0.006]} & 0.53 \\
\bottomrule
\end{tabular}}

\end{table}
